\documentclass{beamer}

\usepackage[utf8]{inputenc}
\usepackage{default}
% Texto Justificado
\usepackage{ragged2e}
\justifying
%  Tema Frankfurt
\usetheme{Frankfurt}

% Slides Negros (Pensar melhor nessa opção, pois pode dificultar a leitura)
%\usecolortheme{structure}

%%%%%%%%%%%%%%%%%%%%%%%%%%%%%%%%%%%%%%%%%%%%%%%%%%%%%%%%%%%%%%%%%%%%%%%%%%%%
% Informações gerais da apresentação
%%%%%%%%%%%%%%%%%%%%%%%%%%%%%%%%%%%%%%%%%%%%%%%%%%%%%%%%%%%%%%%%%%%%%%%%%%%%
\title[TCC]{Implantação e Análise de serviço Web Proxy Cache em Infraestruturas de Computação na Nuvem}
\author{Luís Eduardo Tenório Silva\\\textit{lets@cin.ufpe.br}}
\date[18/12/13]{18 de Dezembro de 2013}
\logo{
    \includegraphics[height=1cm,keepaspectratio]{imagens/ufal.jpg}
    \includegraphics[height=1cm,keepaspectratio]{imagens/comp.png}    
}

% Estruturação:
%    Definição de computação em nuvem
%    Motivação (Computação em nuvem): porque utilizar computação em nuvem.
%    Definição de proxy cache
%    Motivação (proxy cache): porque utilizar proxy cache.
%    Computação em nuvem + proxy cache (Trabalhos realizados) OBS: Pesquisar no IEEE para definir a data onde apareceu a junção dos conceitos.
%    Problemática:
%    Como realiar essa junção? Apresentar em números...OBS: Apresentar agora os problemas das diversas configurações
%    Como deve ser configurado para apresentar melhores resultados? DEPENDE DA MÉTRICA A SER OBSERVADA
%    Proxy cache:
%    - Tipos
%    - Arquiteturas
%    - Protocolos Inter cache
%    - 
%    Cloud Computing:
%    - Características
%    - Categorias
%    - Virtualização
%    Arquiteturas OBS: Descrever cloud+proxy e começar a falar como será realizado a simulação do tráfego
%    Cenários
%    Resultados
%    Considerações finais
%    Justificativa: O porquê escolher cloud + serviço proxy. (Indagar o aumento de serviços na nuvem, E o uso de proxy cache)

% 
\begin{document}
    % Título do trabalho
    \frame{\titlepage}

    \frame{\tableofcontents}

    \section{Computação em nuvem}
    \setbeamercovered{transparent}
    \begin{frame}
      \frametitle{Computação em nuvem}
      \begin{block}{Definição abrangente}<1-| alert@1>
      \textit{Hardwares de grande escala distribuídos geograficamente e infra-estrutura de software composto de recursos heterogêneos de rede compartilhado por várias organizações administrativas que são coordenados para fornecer transparência, confiabilidade, abrangência e suporte computacional persistente para uma grande variedade de aplicações.} (Miguel et
al , 2004)
    \end{block}
    
    \begin{block}{Definição menos abrangente}<2- | alert@2>
      \textit{Forma de computação distribuída em que um "super computador virtual", composto de um conjunto de computadores de baixo
acoplamento interligados em rede, agindo em conjunto para realizar grandes tarefas.} (Buyya et al. 2009)
     \end{block}
       
    \end{frame}
    \setbeamercovered{invisible}

    \section{Motivação (Computação em nuvem)}
    \begin{frame}
      \frametitle{Motivação}
      Por que utilizar computação em nuvem?
      \begin{itemize}
       \item <2->Diversos serviços migrando para a internet e para a nuvem (Email, Editores de texto, Ambientes de desenvolvimento etc).
       \item <3->Economia de recursos
       \item <4->Controle maior do serviço (Elasticidade, pool de recursos...)
      \end{itemize}
    \end{frame}



\end{document}
